\textbf{Пример 1.} При $N = 27$: $3^3 = 27 \le 27 < 64 = 4^3$, поэтому целая часть кубического корня равна $3$. Также $5^2 = 25 \le 27 < 36 = 6^2$, поэтому целая часть квадратного корня, округлённая вниз, равна $5$. Ответ: \texttt{3 5}.

\textbf{Пример 2.} При $N = 63$: $3^3 = 27 \le 63 < 64 = 4^3$, поэтому целая часть кубического корня равна $3$. Также $7^2 = 49 \le 63 < 64 = 8^2$, поэтому целая часть квадратного корня равна $7$. Ответ: \texttt{3 7}. Число $63$ на единицу меньше точной степени ($64 = 4^3 = 8^2$), поэтому оба корня при округлении вниз дают значения на единицу меньше соответствующих показателей степени $4$ и $8$.
